% Generated draft: requires mathematical and source review.
\documentclass[11pt]{article}
\usepackage[margin=25mm]{geometry}
\usepackage{fontspec}
\setmainfont{Latin Modern Roman}
\usepackage{amsmath,amssymb,mathtools,unicode-math}
\setmathfont{Latin Modern Math}
\usepackage{enumitem,needspace,xurl,hyperref}
\setlength{\parindent}{0pt}
\setlength{\parskip}{5pt}
\setlength{\emergencystretch}{3em}
\newcommand{\N}{\mathbb{N}}
\newcommand{\Z}{\mathbb{Z}}
\newcommand{\Q}{\mathbb{Q}}
\newcommand{\R}{\mathbb{R}}
\newcommand{\C}{\mathbb{C}}
\newcommand{\F}{\mathbb{F}}
\newcommand{\A}{\mathbb{A}}
\newcommand{\PP}{\mathbb{P}}
\newcommand{\E}{\mathbb{E}}
\newcommand{\eps}{\varepsilon}
\newcommand{\dd}{\,\mathrm{d}}
\begin{document}
\section*{UnsolvedMath ID: 30000743}
\textbf{Source problem code:} catalogue entry
\subsection{Definitions and mathematical statement}
Let \(n\geq1\), and let \(H=\{A_1,\ldots,A_m\}\) be an \(n\)-uniform hypergraph with pairwise disjoint edges: \(|A_i|=n\) and \(A_i\cap A_j=\varnothing\) for \(i\ne j\). Here \(|H|=m\) counts hyperedges, while the board \(\bigcup_i A_i\) has \(mn\) elements. The source refers first to \(m=2^n\) and asks about the changed regime \(m<2^n\), and about a version in which Picker wants only ``some surplus.'' In the customary pair-offering mechanism, Picker presents two unclaimed board elements, Chooser takes one, and Picker takes the other. The supplied record does not define the surplus, the precise winning objective, or what happens to a last unpaired element. Thus the Picker question is a request to determine the outcome or achievable surplus for these disjoint-edge boards, subject to game rules that still require recovery from the report; it is not a quantified numerical surplus assertion.

For the second question, let \(G\) be an \(n\)-vertex graph and let \(k\geq1\). In the usual \((1,k)\) edge-claiming connectivity game, Maker claims one unclaimed edge per turn and Breaker claims \(k\), with Maker seeking a connected spanning subgraph. The source writes \(c(G,k)\) for the number of Maker's final edges when Maker wins, and \(c(G,k)=\infty\) otherwise, and asks whether
\[
 C_n(k):=\min_{|V(G)|=n}c(G,k)
 \quad\text{satisfies}\quad
 \forall n\ \forall k\geq1,\qquad C_n(k)\leq C_n(k+1).
\]
The minimum uses the extended order in which every finite value is less than \(\infty\). The intended range of \(n\), graph class, first player, stopping rule, and strategy optimization defining a single value \(c(G,k)\) are not stated. For example, minimizing the worst-case number of Maker moves until first connectivity over winning strategies would be an additional specification, not supplied by the phrase ``final edges.'' The problem also requests analogous extremal functions and monotonicity questions for other Maker--Breaker winning objectives; no particular list or universal assertion for those objectives is given.
\subsection{Short English statement}
There are two linked requests: determine Picker's outcome or surplus below the disjoint-edge count \(2^n\), and decide whether the minimum winning-edge count for biased Maker connectivity increases with Breaker's bias. The supplied statement leaves the surplus and the optimization behind the winning-edge count unspecified.
\subsection{Sources}
\begin{itemize}
\item Local UnsolvedMath registry, permanent ID \(30000743\), exact source code \texttt{catalogue entry}, proposed year 2007; no individual proposer is supplied.
\item \emph{Mini-Workshop: Positional Games (2008)}, Oberwolfach Reports, DOI 10.4171/owr/2007/20, pp. 1073--1114. \url{https://doi.org/10.4171/owr/2007/20}.
\item The supplied \texttt{corrected\_verified} note reports normalization after checking mathematical content against the official report. The registry's open label and literature assessment checked 2026-08-21 say that the Picker surplus, monotonicity, and related questions are explicitly left open; this is dated supplied triage, not an independent status verification.
\end{itemize}
\end{document}
